\Article{9}
We now consider, as a second example, the representation of the surface of a right cone in a plane. For the equation of the cone we take
\[
xx+yy-kkzz=0
\]
\noindent where we further set
\[
\begin{aligned}
x&=kt\cos u\\
y&=kt\sin u\\
z&=t
\end{aligned}
\]
\noindent and, as before, put \(X=T,Y=U,Z=0\).

The differential equation
\[
\omega=(kk+1)\mathrm{d}t^2+kktt\,\mathrm{d}u^2=0
\]
\noindent here gives the two integrals
\[
\log t\pm i\BareRoot{\frac{kk}{kk+1}}\mathord{.}u=\text{Const.}
\]
\noindent We therefore have the solution
\[
\begin{aligned}
X+iY&=f(\log t+i\BareRoot{\frac{kk}{kk+1}}\mathord{.}u)\\
X-iY&=f'(\log t-i\BareRoot{\frac{kk}{kk+1}}\mathord{.}u)
\end{aligned}
\]
\noindent That is, with \(f\) denoting an arbitrary function, we take for \(X\) the real part of
\[
f(\log t+i\BareRoot{\frac{kk}{kk+1}}\mathord{.}u)
\]
\noindent and for \(Y\) the imaginary part, with the factor \(i\) omitted.
\par\indent\SourceLeaf{208}%
If, for example, we take an exponential quantity for \(f\), namely
\[
f\upsilon=he^{\upsilon}
\]
\noindent where \(h\) is constant and \(e\) denotes the base of the hyperbolic logarithms, we obtain the simplest representation
\[
X=ht\cos\BareRoot{\frac{kk}{kk+1}}\mathord{.}u,\quad
Y=ht\sin\BareRoot{\frac{kk}{kk+1}}\mathord{.}u
\]

Applying the formulas of Article 7 here gives
\[
n=(kk+1)tt,\quad N=1
\]
\noindent and, since \(\varphi\upsilon=\varphi'\upsilon=he^{\upsilon}\),
\[
\varphi(\log t+i\BareRoot{\frac{kk}{kk+1}}\mathord{.}u)\mathord{.}\varphi'(\log t-i\BareRoot{\frac{kk'}{kk+1}}\mathord{.}u)=hhtt
\]
\noindent hence
\[
m=\frac{h}{\BareRoot{(kk+1)}}
\]
\noindent so the ratio is constant. If we also set
\[
h=\BareRoot{(kk+1)}
\]
\noindent the representation becomes a perfect development.

